%% 2i2d-pv-orientation — irradiance, surface, inclinaison et orientation d'un panneau.
%% Vue de profil : le faisceau solaire (solideD) arrive sur un panneau incliné de bêta ;
%% l'angle d'incidence i est mesuré entre le rayon et la normale au panneau.
\documentclass[tikz,border=4pt]{standalone}
\usepackage{liaisons}
\begin{document}
\begin{tikzpicture}[x=1cm,y=1cm,
  trait/.style={line width=0.8pt, line cap=round, line join=round},
  rayon/.style={-{Stealth[length=4.5pt]}, line width=1pt, draw=solideD},
  cote/.style={{Stealth[length=4pt]}-{Stealth[length=4pt]}, line width=0.6pt},
  txt/.style={font=\scriptsize, inner sep=1.5pt}]

%% ---------- sol ----------------------------------------------------------
\draw[trait] (0.55,0) -- (4.55,0);
\foreach \k in {0.65,0.85,...,4.45} { \draw[line width=0.4pt] (\k,0) -- (\k-0.14,-0.14); }

%% ---------- panneau incliné de bêta = 30° --------------------------------
\begin{scope}[shift={(2.4,0.95)}, rotate=30]
  \draw[solideB, line width=0.8pt, fill=solideB!25] (-1.05,-0.06) rectangle (1.05,0.06);
\end{scope}
\draw[trait] (1.491,0) -- (1.491,0.425)  (3.309,0) -- (3.309,1.475);   % supports

%% ---------- angle d'inclinaison bêta -------------------------------------
\draw[line width=0.5pt, dash pattern=on 2pt off 1.6pt] (1.491,0.425) -- (2.65,0.425);
\draw[solideD, line width=0.9pt] (1.491,0.425) ++(0:0.46) arc (0:30:0.46);
\node[txt, text=solideD, fill=white, inner sep=0.8pt] at (2.03,0.545) {$\beta$};

%% ---------- Soleil et faisceau -------------------------------------------
\fill[solideD] (0.65,2.55) circle (0.19);
\foreach \a in {0,45,...,315} { \draw[solideD, line width=0.7pt]
  (0.65,2.55) ++(\a:0.25) -- ++(\a:0.11); }
\draw[rayon] (1.008,1.041) -- (1.643,0.597);
\draw[rayon] (1.408,1.272) -- (2.043,0.828);
\draw[rayon] (1.808,1.503) -- (2.443,1.059);
\draw[rayon] (2.208,1.734) -- (2.843,1.290);
\node[txt, text=solideD, anchor=west] at (2.55,2.42)
  {$E_{\mathrm{irr}} = 1\,000$ W$\cdot$m$^{-2}$};

%% ---------- normale au panneau et angle d'incidence i --------------------
\draw[line width=0.6pt, dash pattern=on 2pt off 1.6pt] (2.508,1.013) -- (2.033,1.836);
\node[txt, anchor=south east, inner sep=1pt] at (2.06,1.86) {normale};
\draw[line width=0.6pt] (2.508,1.013) ++(120:0.46) arc (120:145:0.46);
\node[txt, fill=white, inner sep=0.8pt] at (2.116,1.441) {$i$};

%% ---------- cote de la surface S -----------------------------------------
\draw[cote] (1.631,0.183) -- (3.449,1.233);
\node[txt, fill=white, inner sep=1pt] at (2.54,0.708) {$S$};

%% ---------- puissance électrique produite --------------------------------
\draw[-{Stealth[length=7pt,width=6pt]}, line width=2pt, draw=solideE] (3.45,1.62) -- (4.45,1.62);
\node[txt, text=solideE, anchor=west] at (4.52,1.72)
  {$P_{\mathrm{crête}} = \eta \times E_{\mathrm{irr}} \times S$};
\node[txt, text=solideE, anchor=west] at (4.52,1.40) {courant continu ($=$)};


%% ---------- rose des vents ------------------------------------------------
\draw[trait] (5.95,0.42) circle (0.34);
\foreach \a/\l/\anc in {90/{N $0^\circ$}/south, 0/{E $90^\circ$}/west,
                        270/{S $180^\circ$}/north, 180/{O}/east} {
  \draw[line width=0.6pt] (5.95,0.42) ++(\a:0.34) -- ++(\a:0.1);
  \node[font=\tiny, anchor=\anc, inner sep=1.5pt] at ($(5.95,0.42)+(\a:0.46)$) {\l};}
\fill (5.95,0.42) circle (1pt);

%% ---------- note ----------------------------------------------------------
\node[txt, anchor=north, text=black!65] at (3.0,-0.40)
  {production maximale quand $i = 0$};
\node[txt, anchor=north, text=black!65] at (3.0,-0.74)
  {si $\beta = 0$ (panneau à plat), l'orientation n'a aucun effet};
\end{tikzpicture}
\end{document}
